\section{Related Work}
\label{sec:related_work}

Our work is situated at the intersection of parameter-space model merging, dynamic routing, and non-linear dynamical systems.

\subsection{Static Model Merging and Task Vectors}
Model merging has emerged as a lightweight, training-free alternative to joint multi-task learning. The cornerstone of modern merging is the concept of a \emph{task vector} \cite{ilharco2022editing}, constructed by subtracting the pre-trained base model parameters from the fine-tuned expert parameters ($V = W_{expert} - W_{base}$). These vectors can be added linearly to edit model behaviors. Subsequent works have attempted to refine static merging by addressing redundant parameters and sign conflicts. For example, \emph{TIES-Merging} \cite{yadav2023ties} prunes low-magnitude values and resolves parameter sign conflicts before averaging, while \emph{DARE} \cite{yu2024language} uses random pruning to sparsify task vectors. Other methods like \emph{FoldMerge} use diffeomorphisms to align non-linear parameter spaces. However, all these methods remain static at test-time, applying a single, frozen set of blending coefficients across all inputs. Consequently, they are unable to adapt to sample-level variance and suffer from severe representational interference when task boundaries overlap.

\subsection{Dynamic and Adaptive Model Merging}
To overcome the limitations of static merging, several works have proposed dynamic or adaptive merging mechanisms. \emph{AdaMerging} \cite{yang2023adamerging} optimizes layer-wise coefficients at test-time by minimizing prediction entropy on the incoming test stream (unsupervised test-time adaptation). However, AdaMerging's coefficient search is decoupled from individual input features, making it slow to adapt to heterogeneous batches. 
To achieve sample-wise dynamic merging, input-dependent routers have been explored, which map input pooled representations to routing coefficients.
Recently, \emph{QWS-Merge} proposed projecting input features into a quantum-like phase space and using cosine wave superposition to compute dynamic coefficients.
While dynamic routers improve on-device adaptability, they often suffer from the \emph{Overfitting-Optimizer Paradox} \cite{trial2_submission3} when optimized on small datasets, or rely on simple, flat linear maps that cannot model complex representational boundaries. In contrast, ChaosMerge enforces a highly regularized, chaotic Coupled Map Lattice (CML) structure that guides the trajectory of merging coefficients dynamically without over-parameterization.

\subsection{Non-linear Dynamics and Chaos in Machine Learning}
The connection between deep neural networks and discrete dynamical systems has a rich history. A deep neural network can be conceptualized as a discrete dynamical system where layer transitions correspond to temporal steps of a trajectory \cite{e2017proposal, haber2017stable}. Under this view, training can be framed as optimal control, and network depth corresponds to simulated time. While previous works have leveraged this connection to study training stability or design invertible architectures, the application of non-linear chaos theory to parameter fusion remains unexplored. 
In physics, a \emph{Coupled Map Lattice} (CML) \cite{kaneko1992coupled} is a standard model for studying spatio-temporal chaos, where localized non-linear maps interact via spatial coupling. ChaosMerge is the first to leverage the unique self-organizing properties of CMLs and chaotic attractor trajectories to regularize and govern the dynamic assembly of multi-task parameters.

\subsection{Parameter-Efficient Fine-Tuning (PEFT) and Mixture-of-Experts (MoE)}
Our work also relates to Parameter-Efficient Fine-Tuning (PEFT) and Mixture-of-Experts (MoE) routing, such as LoRA-MoE \cite{dou2023loramoe} and adaptive parameter steering. Standard MoE methods train multiple specialized sub-networks (experts) and route inputs dynamically using unconstrained Softmax routers. While highly expressive, these routers require substantial memory to store expert parameters and are prone to routing instability or collapse in low-resource regimes. ChaosMerge can be seen as a highly regularized, parameter-efficient parameter steering mechanism that fuses expert weights directly in parameter-space, eliminating the need for active multi-expert memory swapping and providing a stable, physically bounded dynamical steering prior.
